\section{Contributions}\label{sec:contributions}

This thesis delivers an end-to-end system rather than an isolated language prototype.
The contributions span language design, compiler engineering, a production web application, automated quality
enforcement, and quantitative evaluation.
They are summarized below in the order a reader encounters them in the preceding chapters.

\paragraph{Motivo as an offline composition language.}
The first contribution is the Motivo source language for deterministic, structural music composition.
It maps musical entities---notes, rests, chords, sequences, patterns, tracks, and voices---to typed programming
constructs rather than to raw MIDI byte construction.
Patterns with typed parameters compress repetitive structure in the source program; the lowering stage materializes the
full event timeline described in Chapters~\ref{ch:language-design},~\ref{ch:compiler-implementation},
and~\ref{ch:evaluation}.
The language targets Standard MIDI Files as a portable artifact, separating compositional intent from playback-time
synthesis or live scheduling.

\paragraph{Compiler architecture and the version~1.0 to version~2.0 evolution.}
The second contribution is a staged native compiler: scanner, parser, semantic analyzer, cursor-based lowerer, and
MIDI backend.
The project-specific core is the lowering model, in which sequential statements advance musical time, explicit offsets
place events without moving the main cursor, voices maintain isolated cursors, and pattern calls return reconstructed
musical values.
Between version~1.0 and version~2.0 the surface language and semantic pipeline were revised deliberately: implicit
\texttt{let}-based typing and arity-only pattern overload gave way to explicit \texttt{var} declarations,
signature-based
pattern resolution, and single-pass semantic analysis with assignability checks.
Chapters~\ref{ch:language-design} and~\ref{ch:compiler-implementation} document this evolution;
Chapter~\ref{ch:evaluation} shows that the stricter model
imposes roughly $0$--$2\%$ end-to-end wall-clock overhead on large synthetic workloads while producing byte-identical
MIDI output.

\paragraph{Motivo Studio, testing, and deployment.}
The third contribution is Motivo Studio, which embeds the compiler in a full-stack product
(Chapter~\ref{ch:studio-editor}).
A Next.js client provides editing, diagnostics, piano-roll visualization, and browser playback; an Express API wraps
the \texttt{motivoc} binary; PostgreSQL backs authenticated file storage.
The monorepo enforces quality at four boundaries---635 CTest cases for the compiler, 31 server tests, 79 client tests,
and paper CI workflows (Chapter~\ref{ch:deployment-ci}, Section~\ref{sec:testing-ci}).
Production deployment on Railway with Cloudflare DNS and TLS at \texttt{motivo-studio.eu} demonstrates that the thesis
work extends beyond a local prototype to a deployable engineering system with private database networking and
CI-gated promotion.

\paragraph{Empirical evaluation.}
The fourth contribution is bounded quantitative evidence.
Synthetic benchmarks expand compact programs into up to 20,000 note events and compile them in about 5--8~ms on a local
Apple M3 Pro Release build.
Per-stage timing confirms that semantic analysis is a small fraction of total cost relative to lowering and MIDI
serialization.
Together with the test inventory, these results support the claim that Motivo is suitable for interactive local
experimentation at the current scale, without overstating musical or user-experience superiority.
